\documentclass[11pt,a4paper]{article}
% Preâmbulo institucional comum — DEE/IFMA
\usepackage[utf8]{inputenc}
\usepackage[T1]{fontenc}
\usepackage[brazil]{babel}
\usepackage{lmodern}
\usepackage{microtype}
\usepackage{geometry}
\usepackage{xcolor}
\usepackage{graphicx}
\usepackage{booktabs}
\usepackage{tabularx}
\usepackage{array}
\usepackage{enumitem}
\usepackage{hyperref}
\usepackage{lastpage}
\usepackage{fancyhdr}
\usepackage{setspace}
\usepackage{titlesec}
\usepackage{tcolorbox}
\usepackage{ragged2e}

\geometry{a4paper,top=2.2cm,bottom=2.2cm,left=2.5cm,right=2.5cm}
\definecolor{azulpetroleo}{HTML}{0F4C5C}
\definecolor{ambar}{HTML}{C8892D}
\definecolor{cinzaclaro}{HTML}{F2F5F7}
\definecolor{cinzaescuro}{HTML}{374151}
\hypersetup{colorlinks=true,linkcolor=azulpetroleo,urlcolor=azulpetroleo,citecolor=azulpetroleo}
\setlength{\parindent}{1.25cm}
\setlength{\parskip}{0.35em}
\onehalfspacing
\titleformat{\section}{\large\bfseries\color{azulpetroleo}}{}{0pt}{}
\titleformat{\subsection}{\normalsize\bfseries\color{azulpetroleo}}{}{0pt}{}
\pagestyle{fancy}
\fancyhf{}
\fancyhead[L]{\small Departamento de Eletroeletrônica — DEE}
\fancyhead[R]{\small IFMA — Campus São Luís/Monte Castelo}
\fancyfoot[C]{\small Página \thepage\ de \pageref{LastPage}}
\renewcommand{\headrulewidth}{0.4pt}
\renewcommand{\footrulewidth}{0.4pt}
\newcommand{\campo}[1]{\textcolor{ambar}{\textbf{[#1]}}}
\newcommand{\instituicao}{Instituto Federal de Educação, Ciência e Tecnologia do Maranhão}
\newcommand{\campus}{Campus São Luís/Monte Castelo}
\newcommand{\departamento}{Departamento de Eletroeletrônica — DEE}
\newcommand{\cidade}{São Luís — MA}
\newcommand{\cabecalhoInstitucional}{%
\begin{center}
{\Large\bfseries\color{azulpetroleo}\instituicao}\\[2mm]
{\large\bfseries\campus}\\[1mm]
{\normalsize\departamento}
\end{center}
\vspace{2mm}
\hrule
\vspace{4mm}
}
\newtcolorbox{destaque}{colback=cinzaclaro,colframe=azulpetroleo,boxrule=0.6pt,arc=2mm}


\geometry{a4paper,top=1.1cm,bottom=1.1cm,left=2.2cm,right=2.2cm}
\setstretch{1.04}
\setlength{\parindent}{0pt}
\setlength{\parskip}{0.20em}
\pagestyle{empty}

% As caixas lado a lado usam tcbraster: o tabularx alinharia pela primeira
% linha de base, deixando os topos desencontrados quando os textos têm
% alturas diferentes. O raster alinha pelo topo e iguala as alturas da linha.
\tcbuselibrary{raster}

\newtcolorbox{resumobox}[1]{
  title={\sffamily\bfseries #1},
  coltitle=azulpetroleo,
  colback=cinzaclaro,
  colframe=azulpetroleo!30,
  boxrule=0.55pt,
  arc=1.5mm,
  left=3mm,right=3mm,top=1.4mm,bottom=1.4mm,
  before upper={\RaggedRight},
  fonttitle=\small
}

\begin{document}

\begin{center}
\scalebox{0.82}{\marcaIFMA}\\[1.5mm]
{\sffamily\small\bfseries\color{cinzaescuro} CAMPUS SÃO LUÍS -- MONTE CASTELO}\\[0.7mm]
{\sffamily\footnotesize\color{cinzamedio} DEPARTAMENTO DE ELETROELETRÔNICA -- DEE}\\[2mm]
{\color{ambar!85}\rule{0.62\textwidth}{0.8pt}}\\[2.5mm]
{\sffamily\LARGE\bfseries\color{azulpetroleo} Resumo Executivo}\\[1mm]
{\sffamily\large\color{cinzaescuro} Programa de Modernização Tecnológica dos Laboratórios do DEE}
\end{center}

\vspace{1mm}
\begin{resumobox}{Propósito}
Modernizar a infraestrutura de ensino e pesquisa do DEE por meio de doações, cooperação técnica, licenças, capacitação e projetos de P\&D com empresas. A iniciativa atende aos cursos técnicos de Eletrotécnica, Eletrônica e Automação e ao curso de Engenharia Elétrica.
\end{resumobox}

\begin{tcbraster}[raster columns=3, raster equal height=rows,
                  raster column skip=2.5mm, raster left skip=0pt, raster right skip=0pt]
\begin{resumobox}{Capacidade}
\centering{\sffamily\Large\bfseries\color{azulpetroleo} 34}\\[-0.5mm]
{\small docentes em macroáreas complementares}
\end{resumobox}
\begin{resumobox}{Prospecção}
\centering{\sffamily\Large\bfseries\color{azulpetroleo} 60}\\[-0.5mm]
{\small empresas classificadas e organizadas}
\end{resumobox}
\begin{resumobox}{Abrangência}
\centering{\sffamily\Large\bfseries\color{azulpetroleo} 14}\\[-0.5mm]
{\small eixos de modernização tecnológica}
\end{resumobox}
\end{tcbraster}

\begin{tcbraster}[raster columns=2, raster equal height=rows,
                  raster column skip=2.5mm, raster left skip=0pt, raster right skip=0pt]
\begin{resumobox}{Competências estratégicas}
\small Automação e IIoT; eletrônica e sistemas embarcados; IA e aprendizado de máquina; sensores inteligentes; telecomunicações; máquinas e acionamentos; proteção e sistemas de energia; renováveis, armazenamento e microredes; cibersegurança OT; robótica e visão computacional.
\end{resumobox}
\begin{resumobox}{Apoios prioritários}
\small Bancadas e instrumentos; CLPs, IHMs e redes; kits embarcados; workstations com GPU; relés e IEDs; plataformas de software; licenças educacionais; atualização de equipamentos; treinamento de servidores e estudantes.
\end{resumobox}
\end{tcbraster}

\begin{resumobox}{Resultados e compromissos}
\small A parceria deverá fortalecer aulas práticas, iniciação científica, TCCs, pesquisas de mestrado e doutorado, extensão e projetos de P\&D. O DEE prevê registro patrimonial, destinação documentada, acompanhamento de uso, manutenção, segurança, relatórios de resultados e reconhecimento institucional compatível com as normas do IFMA.
\end{resumobox}

\begin{resumobox}{Próximo passo}
\small Reunião técnica para apresentar os laboratórios, confirmar aderência entre as necessidades do DEE e o portfólio do parceiro e definir a modalidade de cooperação. Quantitativos discentes, disciplinas atendidas, orçamento e cronograma serão consolidados na ficha institucional antes da formalização de cada proposta.
\end{resumobox}

\vspace{2mm}
\begin{center}
{\sffamily\bfseries\color{azulpetroleo} Francisco dos Santos Viana}\\[-0.2mm]
{\small Chefe do Departamento de Eletroeletrônica -- DEE}\\[-0.2mm]
{\small\href{mailto:francisco.viana@ifma.edu.br}{francisco.viana@ifma.edu.br} \quad | \quad (98) 98816-6263}\\[1mm]
{\footnotesize\color{cinzamedio} IFMA -- Campus São Luís/Monte Castelo \textbullet\ São Luís -- MA \textbullet\ 2026}
\end{center}

\end{document}
